% !TEX program = xelatex
\documentclass[11pt,a4paper]{article}
\usepackage[margin=2.2cm,top=2cm]{geometry}
\usepackage{fontspec}
\usepackage{polyglossia}
\setdefaultlanguage{turkish}
\usepackage{amsmath,amssymb}
\usepackage{graphicx}
\usepackage{booktabs}
\usepackage{enumitem}
\usepackage{multicol}
\usepackage[most]{tcolorbox}
\usepackage{fancyhdr}
\usepackage{titlesec}
\usepackage{xcolor}
\usepackage{needspace}
\usepackage[hidelinks]{hyperref}

\definecolor{anamavi}{RGB}{0,45,98}
\definecolor{cozumyesil}{RGB}{0,110,50}
\definecolor{uyarikirmizi}{RGB}{180,20,20}

\titleformat{\section}{\large\bfseries\color{anamavi}}{\thesection.}{0.5em}{}
\titleformat{\subsection}{\normalsize\bfseries\color{anamavi}}{\thesubsection}{0.5em}{}

\pagestyle{fancy}\fancyhf{}
\lhead{\small MMAT105 Genel Matematik}
\rhead{\small Hafta 5 -- Ders Notu}
\cfoot{\small\thepage}
\renewcommand{\headrulewidth}{0.4pt}

\newtcolorbox{kural}[1][Kural]{colback=anamavi!5,colframe=anamavi,title=#1,fonttitle=\bfseries,breakable,left=4pt,right=4pt,top=2pt,bottom=2pt}
\newtcolorbox{dikkat}[1][Dikkat]{colback=uyarikirmizi!5,colframe=uyarikirmizi,title=#1,fonttitle=\bfseries,breakable,left=4pt,right=4pt,top=2pt,bottom=2pt}
\newcounter{ornek}
\newenvironment{ornek}{\refstepcounter{ornek}\par\medskip\noindent\textbf{\color{anamavi}Örnek \theornek.}\ }{\par}
\newenvironment{cozum}{\begin{tcolorbox}[colback=cozumyesil!6,colframe=cozumyesil,boxrule=0.5pt,left=4pt,right=4pt,top=2pt,bottom=2pt,breakable]\textbf{\color{cozumyesil}Çözüm.}\ }{\end{tcolorbox}}
\newcounter{alis}
\newenvironment{alistirma}{\refstepcounter{alis}\par\medskip\needspace{9\baselineskip}\noindent\textbf{\color{anamavi}Alıştırma \thealis.}\ }{\par}
\newcommand{\kendinizoncesi}{\needspace{12\baselineskip}}
\newcommand{\ab}[1]{\overline{#1}}
\setlength{\parindent}{0pt}
\setlength{\parskip}{3pt}

\begin{document}

% ---------------- Başlık ----------------
\begin{minipage}[c]{\textwidth}
{\LARGE\bfseries\color{anamavi} MMAT105 Genel Matematik}\\[4pt]
{\large Hafta 5 Ders Notu: Birinci Dereceden\\ Bir Bilinmeyenli Denklemler}\\[6pt]
{\small 2026--2027 Güz}
\end{minipage}

\vspace{4pt}\noindent\textcolor{anamavi}{\rule{\textwidth}{1pt}}

\textbf{Bu notun kullanımı.} Her bölümde önce kısa bir kural kutusu, ardından çözümlü örnekler yer alır. Örnekleri çözümüne bakmadan önce kendiniz deneyin. Notun sonunda önce çözümlü alıştırmalar, ardından cevap anahtarlı ``Kendiniz Çözün'' alıştırmaları yer alır.

% =====================================================
\section{Denklem, Eşitlik ve Özdeşlik}

\begin{kural}[Tanımlar]
\begin{itemize}[leftmargin=*,itemsep=1pt]
\item \textbf{Eşitlik:} İki ifadenin eşit olduğunu bildiren yazı (ör. \(3+5=8\)).
\item \textbf{Denklem:} Bilinmeyen içeren ve bilinmeyenin her değeri için doğru \textbf{olmayan} eşitlik (ör. \(2x+3=11\), yalnız \(x=4\) için doğru). Bir denklem bazı değerler için doğru olabilir ya da hiçbir değer için doğru olmayabilir (\(\text{ÇK}=\emptyset\)).
\item \textbf{Özdeşlik:} Bilinmeyenin \textbf{her} değeri için doğru olan eşitlik (ör. \(2(x+3)=2x+6\)).
\item Denklemi sağlayan sayıya \textbf{kök} (çözüm), köklerin oluşturduğu kümeye \textbf{çözüm kümesi} denir ve \(\text{ÇK}\) ile gösterilir.
\end{itemize}
Bir denklem, dengedeki bir \textbf{terazi} gibidir: iki kefeye aynı işlem yapılırsa denge bozulmaz.
\end{kural}

\begin{kural}[{{Eşitliğin özellikleri ve \(ax+b=0\) denklemi}}]
\(a=b\) ise \quad \(a+c=b+c\), \quad \(a-c=b-c\), \quad \(a\cdot c=b\cdot c\), \quad \(c\neq 0\) ise \(\dfrac{a}{c}=\dfrac{b}{c}\).
\[ax+b=0\quad(a\neq 0)\quad\Longrightarrow\quad x=-\frac{b}{a},\qquad \text{ÇK}=\left\{-\frac{b}{a}\right\}\]
\textbf{Adımlar:} parantezleri aç \(\to\) bilinmeyenli terimleri bir tarafa, sayıları diğer tarafa topla \(\to\) \(x\)'in katsayısına böl \(\to\) bulduğun değeri denklemde yerine koyarak \textbf{sağla}.
\end{kural}

\begin{dikkat}[Sık yapılan hatalar]
Karşı tarafa geçen terimin işareti değişir: \(x+5=9\Rightarrow x=9-5\). \quad Parantezin önündeki eksi, içteki \textbf{her} terimin işaretini değiştirir: \(-(x-3)=-x+3\). \quad Bilinmeyen içeren bir ifadeye bölmeyin (\(0\) olabilir); \textbf{sıfıra bölme tanımsızdır}.
\end{dikkat}

\begin{ornek}
Dengedeki bir terazinin sol kefesinde \(4\) özdeş top ve \(3\) kg, sağ kefesinde \(1\) top ve \(15\) kg vardır. Bir top kaç kg'dır?
\end{ornek}
\begin{cozum}
Top \(x\) kg olsun: \(4x+3=x+15\). İki kefeden birer top ve \(3\) kg çıkaralım:
\[3x=12\ \Rightarrow\ x=4\text{ kg}\qquad(\text{sağlama: } 16+3=19=4+15)\]
\end{cozum}

\begin{ornek}
\(7x-3=4x+12\) denkleminin çözüm kümesini bulunuz.
\end{ornek}
\begin{cozum}
\[7x-4x=12+3\ \Rightarrow\ 3x=15\ \Rightarrow\ x=5\qquad \text{ÇK}=\{5\}\]
Sağlama: \(35-3=32\), \ \(20+12=32\).
\end{cozum}

\begin{ornek}
\(-6x+9=0\) denkleminin çözüm kümesini bulunuz.
\end{ornek}
\begin{cozum}
\(a=-6,\ b=9\): \quad \(x=-\dfrac{b}{a}=-\dfrac{9}{-6}=\dfrac{3}{2}\). \quad \(\text{ÇK}=\left\{\dfrac{3}{2}\right\}\). \ Sağlama: \(-6\cdot\dfrac{3}{2}+9=-9+9=0\).
\end{cozum}

% =====================================================
\section{Parantezli Denklemler}

\begin{kural}[Parantez açma]
\(a(b+c)=ab+ac\), \quad \(-(b-c)=-b+c\). İç içe parantezlerde \textbf{en içteki} parantezden başlanır: \( (\ )\rightarrow[\ ]\rightarrow\{\ \}\).
\end{kural}

\begin{ornek}
\(4(2x-3)-3(x+1)=2(x+3)\) denklemini çözünüz.
\end{ornek}
\begin{cozum}
\[8x-12-3x-3=2x+6\ \Rightarrow\ 5x-15=2x+6\ \Rightarrow\ 3x=21\ \Rightarrow\ x=7\]
Sağlama: \(4\cdot 11-3\cdot 8=44-24=20\), \ \(2\cdot 10=20\).
\end{cozum}

\begin{ornek}
\(3-\big[2x-(x-5)\big]=4-2x\) denklemini çözünüz.
\end{ornek}
\begin{cozum}
İçten dışa: \(2x-(x-5)=2x-x+5=x+5\).
\[3-(x+5)=4-2x\ \Rightarrow\ -x-2=4-2x\ \Rightarrow\ x=6\]
Sağlama: \(3-[12-1]=-8\), \ \(4-12=-8\).
\end{cozum}

% =====================================================
\section{Kesirli ve Ondalıklı Katsayılı Denklemler}

\begin{kural}[Paydaları yok etme]
\begin{itemize}[leftmargin=*,itemsep=1pt]
\item \textbf{Kesirli denklem:} Denklemin iki tarafındaki \textbf{her terim} paydaların EKOK'u ile çarpılır.
\item \textbf{Ondalıklı denklem:} Terimlerdeki en fazla ondalık basamak sayısı \(k\) ise her terim \(10^{k}\) ile çarpılır.
\end{itemize}
\end{kural}

\begin{dikkat}
Kesir çizgisi parantez gibidir: \(-\dfrac{x-1}{4}\) terimi \(4\) ile çarpılınca \(-(x-1)=-x+1\) olur. Paydası olmayan terimler de EKOK ile çarpılır.
\end{dikkat}

\begin{ornek}
\(\dfrac{x+1}{2}+\dfrac{x-2}{3}=4\) denklemini çözünüz.
\end{ornek}
\begin{cozum}
\(\text{EKOK}(2,3)=6\):
\[3(x+1)+2(x-2)=24\ \Rightarrow\ 3x+3+2x-4=24\ \Rightarrow\ 5x=25\ \Rightarrow\ x=5\]
Sağlama: \(\dfrac{6}{2}+\dfrac{3}{3}=3+1=4\).
\end{cozum}

\begin{ornek}
\(x-\dfrac{x-1}{4}=\dfrac{x+5}{2}\) denklemini çözünüz.
\end{ornek}
\begin{cozum}
\(\text{EKOK}(4,2)=4\); paydasız \(x\) terimi de çarpılır:
\[4x-(x-1)=2(x+5)\ \Rightarrow\ 3x+1=2x+10\ \Rightarrow\ x=9\]
Sağlama: \(9-\dfrac{8}{4}=7\), \ \(\dfrac{14}{2}=7\).
\end{cozum}

\begin{ornek}
\(0{,}25x+0{,}5=0{,}2x+0{,}75\) denklemini çözünüz.
\end{ornek}
\begin{cozum}
En fazla \(2\) ondalık basamak var; her terimi \(100\) ile çarpalım:
\[25x+50=20x+75\ \Rightarrow\ 5x=25\ \Rightarrow\ x=5\]
Sağlama: \(1{,}25+0{,}5=1{,}75\), \ \(1+0{,}75=1{,}75\).
\end{cozum}

% =====================================================
\section{Özel Durumlar ve Parametreli Denklemler}

\begin{kural}[{{\(ax+b=0\) denkleminin çözüm kümesi}}]
\begin{center}
\renewcommand{\arraystretch}{1.3}
\begin{tabular}{@{}lll@{}}
\toprule
Durum & Anlamı & Çözüm kümesi\\\midrule
\(a\neq 0\) & tek kök & \(\text{ÇK}=\left\{-\frac{b}{a}\right\}\)\\
\(a=0\) ve \(b=0\) & \(0\cdot x=0\): her gerçek sayı sağlar (özdeşlik) & \(\text{ÇK}=\mathbb{R}\)\\
\(a=0\) ve \(b\neq 0\) & \(0\cdot x=-b\neq 0\): hiçbir sayı sağlamaz & \(\text{ÇK}=\emptyset\)\\\bottomrule
\end{tabular}
\end{center}
\(ax+b=cx+d\) denkleminde: \ \(\text{ÇK}=\mathbb{R}\Leftrightarrow a=c\) ve \(b=d\); \quad \(\text{ÇK}=\emptyset\Leftrightarrow a=c\) ve \(b\neq d\); \quad \(a\neq c\) ise tek kök vardır.\\
Bir sayı denklemin kökü ise, \(x\) yerine yazıldığında denklemi sağlar.
\end{kural}

\begin{dikkat}
\(\emptyset\) (boş küme) ile \(\{0\}\) farklıdır: \(3x=0\) denkleminin çözüm kümesi \(\{0\}\)'dır, boş küme değildir.
\end{dikkat}

\begin{ornek}
Çözüm kümelerini bulunuz: \ (a) \(4(x-1)+2=2(2x-1)\) \qquad (b) \(5x-3(x+2)=2x+1\)
\end{ornek}
\begin{cozum}
(a) \(4x-2=4x-2\ \Rightarrow\ 0\cdot x=0\): her \(x\) sağlar. \ \(\text{ÇK}=\mathbb{R}\).\\
(b) \(2x-6=2x+1\ \Rightarrow\ 0\cdot x=7\): hiçbir \(x\) sağlamaz. \ \(\text{ÇK}=\emptyset\).
\end{cozum}

\begin{ornek}
\((2a-1)x+3=5x+b-1\) denkleminin çözüm kümesi \(\mathbb{R}\) ise \(a\cdot b\) kaçtır?
\end{ornek}
\begin{cozum}
Katsayılar ve sabitler eşit olmalı: \(2a-1=5\Rightarrow a=3\); \quad \(3=b-1\Rightarrow b=4\). \quad \(a\cdot b=12\).
\end{cozum}

\begin{ornek}
\((k+2)x-5=4x+1\) denklemi veriliyor. (a) Çözüm kümesi boş küme ise \(k\) kaçtır? (b) Denklemin kökü \(3\) ise \(k\) kaçtır?
\end{ornek}
\begin{cozum}
(a) \(k+2=4\Rightarrow k=2\); sabitler \(-5\neq 1\) olduğundan gerçekten \(\text{ÇK}=\emptyset\).\\
(b) \(x=3\) yazılır: \(3(k+2)-5=12+1\Rightarrow 3k+1=13\Rightarrow k=4\). \ Kontrol: \(6x-5=4x+1\Rightarrow x=3\).
\end{cozum}

\begin{ornek}
\(\dfrac{x+m}{3}-\dfrac{x-m}{2}=1\) denkleminin kökü \(4\) ise \(m\) kaçtır?
\end{ornek}
\begin{cozum}
\(x=4\) yazalım ve \(6\) ile çarpalım:
\[2(4+m)-3(4-m)=6\ \Rightarrow\ 8+2m-12+3m=6\ \Rightarrow\ 5m=10\ \Rightarrow\ m=2\]
Kontrol: \(\dfrac{6}{3}-\dfrac{2}{2}=2-1=1\).
\end{cozum}

% =====================================================
\section{Paydada Bilinmeyen Olan Denklemler}

\begin{kural}[Çözüm yolu]
\begin{enumerate}[leftmargin=*,itemsep=1pt]
\item Paydayı sıfır yapan değerleri yaz ve dışarıda bırak (\textbf{tanımsızlık koşulu}).
\item Her terimi paydaların EKOK'u ile çarparak paydaları yok et; denklemi çöz.
\item Bulunan değer tanımsızlık koşuluna uymuyorsa kök değildir, çözüm kümesine alınmaz.
\end{enumerate}
\end{kural}

\begin{ornek}
\(\dfrac{3x-1}{x+2}=2\) denklemini çözünüz.
\end{ornek}
\begin{cozum}
Koşul: \(x\neq-2\). \quad \(3x-1=2(x+2)=2x+4\Rightarrow x=5\). \quad Sağlama: \(\dfrac{14}{7}=2\). \quad \(\text{ÇK}=\{5\}\).
\end{cozum}

\begin{ornek}
\(\dfrac{5}{x}-\dfrac{2}{3x}=\dfrac{13}{6}\) denklemini çözünüz.
\end{ornek}
\begin{cozum}
Koşul: \(x\neq 0\). \ \(\text{EKOK}(x,3x,6)=6x\): \quad \(30-4=13x\Rightarrow x=2\).\\
Sağlama: \(\dfrac{5}{2}-\dfrac{2}{6}=\dfrac{15}{6}-\dfrac{2}{6}=\dfrac{13}{6}\). \quad \(\text{ÇK}=\{2\}\).
\end{cozum}

\begin{ornek}
\(\dfrac{3}{x-2}+1=\dfrac{x+1}{x-2}\) denkleminin çözüm kümesini bulunuz.
\end{ornek}
\begin{cozum}
Koşul: \(x\neq 2\). Her terimi \((x-2)\) ile çarpalım:
\[3+(x-2)=x+1\ \Rightarrow\ x+1=x+1\ \Rightarrow\ 0\cdot x=0\]
Eşitlik her \(x\) için doğrudur; ancak \(x=2\) tanımsızdır. \quad \(\text{ÇK}=\mathbb{R}\setminus\{2\}\).
\end{cozum}

% =====================================================
\section{Sözel İfadeler ve Problemler}

\begin{kural}[Sözel ifadeden denkleme]
\begin{center}
\renewcommand{\arraystretch}{1.25}
\begin{tabular}{@{}ll@{\qquad}ll@{}}
\toprule
İfade (sayı \(x\)) & Cebirsel & İfade & Cebirsel\\\midrule
\(3\) katının \(5\) fazlası & \(3x+5\) & \(5\) fazlasının \(3\) katı & \(3(x+5)\)\\
yarısının \(4\) eksiği & \(\frac{x}{2}-4\) & \(4\) eksiğinin yarısı & \(\frac{x-4}{2}\)\\
ardışık tam sayılar & \(x,\ x+1,\ x+2,\dots\) & ardışık çift / tek sayılar & \(x,\ x+2,\ x+4,\dots\)\\\bottomrule
\end{tabular}
\end{center}
\textbf{Adımlar:} bilinmeyeni seç \(\to\) diğer nicelikleri \(x\) ile yaz \(\to\) denklemi kur ve çöz \(\to\) sağla ve \textbf{sorulan niceliği} yaz.\\
\textbf{Yaş:} \(t\) yıl sonra herkesin yaşı \(t\) artar. \quad \textbf{Yol:} yol \(=\) hız \(\times\) zaman.
\end{kural}

\begin{dikkat}
``sayının \(5\) eksiği'' \(x-5\) iken ``\(5\)'in sayı kadar eksiği'' \(5-x\)'tir. Son adımda sorulan her zaman \(x\) değildir (en büyük sayı, toplam, diğer kişinin yaşı\dots).
\end{dikkat}

\begin{ornek}
Bir sayının \(2\) katının \(3\) fazlası, aynı sayının \(4\) katının \(9\) eksiğine eşittir. Bu sayı kaçtır?
\end{ornek}
\begin{cozum}
\(2x+3=4x-9\Rightarrow 12=2x\Rightarrow x=6\). \quad Sağlama: \(15=24-9\).
\end{cozum}

\begin{ornek}
Ardışık beş tam sayının toplamı \(115\) ise en küçüğü kaçtır?
\end{ornek}
\begin{cozum}
Sayılar \(x,\ x+1,\ x+2,\ x+3,\ x+4\): \quad \(5x+10=115\Rightarrow x=21\). \quad (Sayılar \(21,\dots,25\).)
\end{cozum}

\begin{ornek}
\(40\) kişilik bir sınıfta kız öğrenci sayısı erkek öğrenci sayısından \(6\) fazladır. Sınıfta kaç kız öğrenci vardır?
\end{ornek}
\begin{cozum}
Erkek \(x\), kız \(x+6\): \quad \(2x+6=40\Rightarrow x=17\). \quad Kız öğrenci sayısı \(17+6=23\).
\end{cozum}

\begin{ornek}
Ayşe \(10\), annesi \(38\) yaşındadır. Kaç yıl önce annenin yaşı Ayşe'nin yaşının \(5\) katıydı?
\end{ornek}
\begin{cozum}
\(t\) yıl önce: \(38-t=5(10-t)\Rightarrow 38-t=50-5t\Rightarrow 4t=12\Rightarrow t=3\). \quad Sağlama: \(35=5\cdot 7\).
\end{cozum}

\begin{ornek}
Bir araç A şehrinden B şehrine saatte \(60\) km hızla gidip saatte \(90\) km hızla dönüyor. Gidiş--dönüş toplam \(5\) saat sürdüğüne göre A ile B arası kaç km'dir?
\end{ornek}
\begin{cozum}
Uzaklık \(x\) km; zaman \(=\) yol \(\div\) hız: \quad \(\dfrac{x}{60}+\dfrac{x}{90}=5\). \ \(\text{EKOK}(60,90)=180\):
\[3x+2x=900\ \Rightarrow\ x=180\text{ km}\qquad(\text{sağlama: } 3+2=5\text{ saat})\]
\end{cozum}

\begin{ornek}
Bir cüzdanda yalnızca \(5\) TL'lik ve \(10\) TL'lik banknotlar vardır. Toplam \(20\) banknot ve \(145\) TL olduğuna göre kaç tane \(10\) TL'lik banknot vardır?
\end{ornek}
\begin{cozum}
\(10\) TL'lik \(x\) tane ise \(5\) TL'lik \(20-x\) tanedir: \quad \(10x+5(20-x)=145\Rightarrow 5x+100=145\Rightarrow x=9\). \ Sağlama: \(9\cdot 10+11\cdot 5=90+55=145\).
\end{cozum}

% =====================================================
\section{Özet: Formül Kartı}
\begin{center}
\small
\renewcommand{\arraystretch}{1.18}
\begin{tabular}{@{}lp{0.64\textwidth}@{}}
\toprule
Konu & Formül / Kural\\\midrule
Temel denklem & \(ax+b=0,\ a\neq 0\ \Rightarrow\ x=-\dfrac{b}{a}\)\\
Eşitliğin özellikleri & iki tarafa aynı sayı eklenir/çıkarılır; aynı sayıyla çarpılır; \(0\)'dan farklı sayıya bölünür\\
Kesirli denklem & her terim paydaların EKOK'u ile çarpılır\\
Ondalıklı denklem & \(k\) ondalık basamak \(\Rightarrow\) her terim \(10^{k}\) ile çarpılır\\
\(\text{ÇK}=\mathbb{R}\) & \(ax+b=cx+d\): \ \(a=c\) ve \(b=d\) (özdeşlik)\\
\(\text{ÇK}=\emptyset\) & \(ax+b=cx+d\): \ \(a=c\) ve \(b\neq d\)\\
Kökü verilen denklem & kök \(x\) yerine yazılır, parametre bulunur\\
Paydada bilinmeyen & önce payda \(\neq 0\) koşulu; koşulu bozan değer atılır\\
Ardışık sayılar & \(x,\ x+1,\ x+2\); \ çift/tek: \(x,\ x+2,\ x+4\)\\
Yaş / yol & \(t\) yıl sonra herkes \(+t\); \ yol \(=\) hız \(\times\) zaman\\\bottomrule
\end{tabular}
\end{center}

% =====================================================
\newpage
% Hafta 5 -- alıştırma gövdesi (ders notu ve ayrı alıştırma kağıdı tarafından \input edilir)
\section{Çözümlü Alıştırmalar}
Her soruyu önce kendiniz deneyin, sonra çözümü okuyun.

\begin{alistirma}
\(2\big[3-(x-1)\big]=x+2\) denklemini çözünüz.
\end{alistirma}
\begin{cozum}
En içteki parantezden başlanır: \(3-(x-1)=3-x+1=4-x\).\\
\(2(4-x)=x+2\ \Rightarrow\ 8-2x=x+2\ \Rightarrow\ 8-2=x+2x\ \Rightarrow\ 6=3x\ \Rightarrow\ x=2\).\\
Sağlama: \(2\big[3-1\big]=4\), \ \(2+2=4\) \checkmark \qquad \(\text{ÇK}=\{2\}\)
\end{cozum}

\begin{alistirma}
\(\dfrac{x+2}{3}-\dfrac{x-2}{5}=2\) denklemini çözünüz.
\end{alistirma}
\begin{cozum}
Her terim \(\text{EKOK}(3,5)=15\) ile çarpılır (sağdaki \(2\) dahil):
\[5(x+2)-3(x-2)=30\ \Rightarrow\ 5x+10-3x+6=30\ \Rightarrow\ 2x+16=30\ \Rightarrow\ x=7\]
Dikkat: \(-3(x-2)=-3x+6\). \quad Sağlama: \(\dfrac{9}{3}-\dfrac{5}{5}=3-1=2\) \checkmark \qquad \(\text{ÇK}=\{7\}\)
\end{cozum}

\begin{alistirma}
\(0{,}15x+0{,}4=0{,}05x+1{,}4\) denklemini çözünüz.
\end{alistirma}
\begin{cozum}
En çok iki ondalık basamak var: her terim \(100\) ile çarpılır:
\[15x+40=5x+140\ \Rightarrow\ 10x=100\ \Rightarrow\ x=10\]
Sağlama: \(1{,}5+0{,}4=1{,}9\), \ \(0{,}5+1{,}4=1{,}9\) \checkmark \qquad \(\text{ÇK}=\{10\}\)
\end{cozum}

\begin{alistirma}
\((2m+1)x-7=5x+n-2\) denkleminin çözüm kümesi \(\mathbb{R}\) ise \(m\cdot n\) kaçtır?
\end{alistirma}
\begin{cozum}
Çözüm kümesi \(\mathbb{R}\) ise denklem bir özdeşliktir: \(x\)'in katsayıları ve sabit terimler eşit olmalıdır.\\
\(x\)'in katsayıları: \(2m+1=5\Rightarrow m=2\). \quad Sabitler: \(-7=n-2\Rightarrow n=-5\).\\
\(m\cdot n=2\cdot(-5)=-10\). \ (Kontrol: \(5x-7=5x-7\) her \(x\) için doğrudur.)
\end{cozum}

\begin{alistirma}
\(\dfrac{2x}{x-4}=\dfrac{8}{x-4}\) denkleminin çözüm kümesini bulunuz.
\end{alistirma}
\begin{cozum}
Koşul: \(x-4\neq 0\Rightarrow x\neq 4\). \ Paydalar eşit ve sıfırdan farklı olduğundan paylar eşit olmalı: \(2x=8\Rightarrow x=4\).\\
Bulunan değer koşulu bozar (\(x=4\) için paydalar \(0\) olur), atılır. \quad \(\text{ÇK}=\emptyset\).
\end{cozum}

\begin{alistirma}
\(\dfrac{4}{x}+\dfrac{1}{2}=\dfrac{7}{2x}\) denkleminin çözüm kümesini bulunuz.
\end{alistirma}
\begin{cozum}
Koşul: \(x\neq 0\). \ Paydaların EKOK'u \(2x\); her terim \(2x\) ile çarpılır:
\[8+x=7\ \Rightarrow\ x=-1\qquad(-1\neq 0\ \checkmark)\]
Sağlama: \(\dfrac{4}{-1}+\dfrac{1}{2}=-\dfrac{7}{2}\), \ \(\dfrac{7}{-2}=-\dfrac{7}{2}\) \checkmark \qquad \(\text{ÇK}=\{-1\}\)
\end{cozum}

\begin{alistirma}
Bir sayının \(5\) fazlasının \(2\) katı, aynı sayının \(3\) katından \(4\) eksiktir. Bu sayı kaçtır?
\end{alistirma}
\begin{cozum}
Sayı \(x\): \ ``\(5\) fazlasının \(2\) katı'' \(=2(x+5)\); \ ``\(3\) katından \(4\) eksik'' \(=3x-4\).
\[2(x+5)=3x-4\ \Rightarrow\ 2x+10=3x-4\ \Rightarrow\ x=14\]
Sağlama: \(2\cdot 19=38\), \ \(3\cdot 14-4=38\) \checkmark
\end{cozum}

\begin{alistirma}
Aralarında \(390\) km olan iki şehirden aynı anda birbirine doğru yola çıkan iki araçtan biri diğerinden saatte \(10\) km daha hızlı gitmektedir. Araçlar \(3\) saat sonra karşılaştığına göre yavaş olan aracın hızı saatte kaç km'dir?
\end{alistirma}
\begin{cozum}
Yavaş aracın hızı \(x\) km/sa, hızlı olanınki \(x+10\) km/sa. Karşılaşana kadar iki aracın aldığı yolların toplamı \(390\) km'dir (yol \(=\) hız \(\times\) zaman):
\[3x+3(x+10)=390\ \Rightarrow\ 6x+30=390\ \Rightarrow\ 6x=360\ \Rightarrow\ x=60\]
Sağlama: \(3\cdot 60+3\cdot 70=180+210=390\) \checkmark \quad Yavaş aracın hızı \(60\) km/sa.
\end{cozum}

\kendinizoncesi
\section{Alıştırmalar (Kendiniz Çözün)}

\begin{enumerate}[label=\textbf{\arabic*.},leftmargin=*,itemsep=6pt]
\item \(6x-5=2x+19\) denklemini çözünüz.
\item \(3(x-4)+2=5-(x-1)\) denklemini çözünüz.
\item \(\dfrac{x}{4}-\dfrac{x}{6}=2\) denklemini çözünüz.
\item \(0{,}4x-1{,}2=0{,}1x+0{,}3\) denklemini çözünüz.
\item Çözüm kümelerini bulunuz: \ (a) \(3(x+2)-x=2(x+3)\) \qquad (b) \(4x-2(2x-1)=5\)
\item \(3x+2k=kx+11\) denkleminin kökü \(3\) ise \(k\) kaçtır?
\item \((a-4)x+6=2x+3\) denkleminin çözüm kümesi boş küme ise \(a\) kaçtır?
\item \(\dfrac{x+7}{x-1}=3\) denkleminin çözüm kümesini bulunuz.
\item Ardışık üç tek sayının toplamı \(57\) ise en büyüğü kaçtır?
\item Bir baba \(45\), oğlu \(15\) yaşındadır. Kaç yıl sonra babanın yaşı oğlunun yaşının \(2\) katı olur?
\item \(\dfrac{x-1}{2}+\dfrac{x+1}{3}=x-2\) denklemini çözünüz.
\item Bir sınıftaki öğrencilerin \(\frac{3}{5}\)'i kızdır. Sınıfa \(4\) erkek öğrenci katılınca kızların sayısı erkeklerin sayısına eşit oluyor. Başlangıçta sınıfta kaç öğrenci vardır?
\end{enumerate}

\subsection*{Cevap Anahtarı}
{\small
\begin{multicols}{2}
\begin{enumerate}[label=\textbf{\arabic*.},leftmargin=*,itemsep=1pt]
\item \(4x=24\Rightarrow x=6\).
\item \(3x-10=6-x\Rightarrow 4x=16\Rightarrow x=4\).
\item \(\cdot 12\): \(3x-2x=24\Rightarrow x=24\).
\item \(\cdot 10\): \(4x-12=x+3\Rightarrow x=5\).
\item (a) \(0\cdot x=0\): \(\text{ÇK}=\mathbb{R}\). \ (b) \(0\cdot x=3\): \(\text{ÇK}=\emptyset\).
\item \(x=3\): \(9+2k=3k+11\Rightarrow k=-2\).
\item \(a-4=2\Rightarrow a=6\) \ (sabitler: \(6\neq 3\)).
\item \(x\neq 1\); \ \(x+7=3x-3\Rightarrow \text{ÇK}=\{5\}\).
\item \(3x+6=57\Rightarrow x=17\); \ \(17,19,21\): \(21\).
\item \(45+t=2(15+t)\Rightarrow t=15\) yıl sonra.
\item \(\cdot 6\): \(3x-3+2x+2=6x-12\Rightarrow x=11\).
\item Kız \(\frac{3x}{5}\), erkek \(\frac{2x}{5}+4\); eşitlik: \(\frac{x}{5}=4\Rightarrow x=20\).
\end{enumerate}
\end{multicols}}


\vspace{6pt}
\noindent\textcolor{anamavi}{\rule{\textwidth}{0.5pt}}\\
{\small Gelecek hafta: Birinci dereceden iki bilinmeyenli denklemler.}

\end{document}
